\documentclass[10pt,a4paper]{amsart}
\usepackage[margin=19mm]{geometry}
\usepackage{amsmath,amssymb,mathtools}
\usepackage[scr=rsfs,cal=euler]{mathalfa}
\usepackage{xfrac}
\usepackage[unicode,hidelinks,bookmarks=false]{hyperref}
\newcommand{\ie}{i.e.\ }
\newcommand{\cf}{cf.\ }
\newcommand{\resp}{respectively\ }
\newcommand{\Subsection}{\subsection}
\providecommand{\nobreakdash}{\nobreak\hbox{-}}
\setlength{\emergencystretch}{2em}
\multlinegap0pt
\usepackage{bm}
\usepackage{bbm}
\usepackage{dsfont}
\usepackage{xspace}
\usepackage{cancel}
\usepackage{mathtools}
\usepackage{amsthm}
\makeatletter
\@ifundefined{theorem}{\newtheorem{theorem}{Theorem}}{}
\@ifundefined{lemma}{\newtheorem{lemma}[theorem]{Lemma}}{}
\@ifundefined{proposition}{\newtheorem{proposition}[theorem]{Proposition}}{}
\makeatother
\title[Weakly Strongly 2-Nil-Clean Rings]{Weakly Strongly 2-Nil-Clean Rings}
\author{Peter Danchev, Mina Doostalizadeh,, Ahmad Moussavi}
\date{Source: arXiv:2509.05673v1 (CC BY 4.0)}
\begin{document}
\maketitle

\noindent\emph{Source and licence.} Weakly Strongly 2-Nil-Clean Rings. Version arXiv:2509.05673v1 (CC BY 4.0), distributed under the Creative Commons Attribution 4.0 International licence, as recorded in the arXiv licence metadata for this exact version and shown on the abstract page https://arxiv.org/abs/2509.05673v1. Full rights notice: licenses/arxiv-2509.05673-CC-BY-4.0.txt.\par\smallskip

\noindent\textit{Editorial scope.} Counted unit: the Proposition whose source label is \texttt{2.2} in the article (source0.tex lines 324--326, source label \texttt{2.2}), delivered together with the article's own complete original proof of it (source0.tex lines 328--345, one \texttt{proof} environment of 18 lines). No mathematics is written, repaired or re-derived: statement and proof are the article's own text line for line. Required context retained uncounted: hom (lines 270--272); 2 (lines 281--282); 3 (lines 289--290); 5 (lines 307--308); 2nil (lines 315--316). Every definition, display and label of those lines is retained unchanged. Omitted from this package and not redistributed: 1--269, 273--280, 283--288, 291--306, 309--314, 317--323, 327--327, 346--683. Nothing inside a retained excerpt is omitted or altered: every retained body is delivered exactly as the article's lines are written, so no omission receipt of any kind accompanies this package. Citations: the retained excerpts carry no citation command at all, so no bibliography entry of the article is transcribed and no reference is redistributed. Reference adaptation: none was needed and none was made; every reference inside the delivered unit resolves inside it, so no unresolved reference or citation is produced. Printed slips kept literal: no printed word, symbol, hypothesis or number of the counted statement or of its proof was altered. Layout and class: the article's own class is \texttt{amsart} with packages including amssymb, amsmath, amscd, amsthm, comment, xcolor, ulem, tikz, float, graphicx, circuitikz, hyperref; the wrapper is the lane's pinned \texttt{amsart} editorial wrapper at 10pt on A4, which restates the theorem-like environments the retained text needs and adds the article's own macro definitions that the retained text needs (none), with extra wrapper packages (bm, bbm, dsfont, xspace, cancel, mathtools). The delivered numbering is the unit's own. Assets: no retained span contains a figure environment, a graphics-inclusion command, a PDF-page inclusion, a table or a TikZ picture; no image file is copied into this package, the package declares no asset dependency and no image file is redistributed with it. Licence: version arXiv:2509.05673v1 of this article is distributed under the Creative Commons Attribution 4.0 International licence, as recorded in the arXiv licence metadata for this exact version and shown on the abstract page https://arxiv.org/abs/2509.05673v1. A full rights notice is delivered at licenses/arxiv-2509.05673-CC-BY-4.0.txt. Characters: every retained excerpt is pure ASCII, so the delivered engine, XeLaTeX with its default Latin Modern text font, reports no missing character.
\begin{proposition}\label{hom} {\rm (i)} Any homomorphic image of a weakly strongly 2-nil-clean ring is again a weakly strongly 2-nil-clean ring, but the reciprocal claim is not generally true. In particular, $R$ is a weakly strongly 2-nil-clean ring if, and only if, the quotient $R/I$ is a weakly strongly 2-nil-clean ring, whenever $I$ is a nil-ideal of $R$.\\\indent
{\rm (ii)} The class of weakly strongly 2-nil-clean rings is not closed under formation of finite products; e.g., $\mathbb{Z}_5 \times \mathbb{Z}_5$ is not weakly strongly 2-nil-clean.
 \end{proposition}

\begin{lemma}\label{2} Let $R$ be a weakly strongly 2-nil-clean ring. If $2\in nil(R)$, then $R$ is strongly nil-clean.
\end{lemma}

\begin{lemma}\label{3} Let $R$ be a weakly strongly 2-nil-clean ring. If $3\in nil(R)$, then $R$ is strongly 2-nil-clean.
\end{lemma}

\begin{lemma}\label{5} Let $R$ be a weakly strongly 2-nil-clean ring. If there are $e,f\in Id(R)$ and $n\in nil(R)$ commuting with each other such that  $-2=e\pm f+n$, then $R$ is strongly 2-nil-clean.
\end{lemma}

\begin{lemma}\label{2nil} Let $R$ be a ring such that $3\in nil(R)$. Then, the following three statements are equivalent:\\\indent (1) $R$ is a strongly 2-nil-clean ring;\\\indent (2) For any $a \in R$, there exist two idempotents $e, f \in R$ and a nilpotent $w\in  R$ which commute with each other such that $a = e - f + w$;\\\indent  (3) For any $a \in R$, there exist two idempotents $e, f \in R$ and a nilpotent $w\in  R$ which commute with each other such that $a = -e - f + w$.
\end{lemma}

\begin{proposition}\label{2.2} Let $\{R_{\alpha}\}$ be a finite collection of rings. Then, the direct product
$\prod_{\alpha} R_{\alpha}$ is weakly strongly 2-nil-clean if, and only if, each $R_{\alpha}$ is weakly strongly 2-nil-clean, and at most one $R_{\alpha}$ is not strongly 2-nil-clean.
\end{proposition}

\begin{proof} Let $\prod_{\alpha} R_{\alpha}$ be weakly 2-nil-clean. Then, each $R_{\alpha}$ is weakly 2-nil-clean, because, in view of Proposition~\ref{hom}, homomorphic images of weakly strongly 2-nil-clean rings are again weakly strongly 2-nil-clean.

Suppose now that there are two indices, say $\alpha_1$ and $\alpha_2$, such that $\alpha_1\neq \alpha_2$ and both $R_{\alpha_1}$ and $R_{\alpha_2}$ are {\it not} 2-nil-clean. But, as $R_{\alpha_1}\times R_{\alpha_2}$ is a homomorphic image of $R = \prod R_{\alpha}$, we can get $R_{\alpha_1}\times R_{\alpha_2}$ is weakly strongly 2-nil-clean. Considering the element $(-2,2)$ of $R_{\alpha_1}\times R_{\alpha_2}$, one finds that there are $(e,f),(g,h)\in Id(R_{\alpha_1}\times R_{\alpha_2})$ and $(n_1,n_2)\in nil(R_{\alpha_1}\times R_{\alpha_2})$ commuting with each other such that $(-2,2)=\pm(e,f)\pm (g,h)+(n_1,n_2)$. If $(-2,2)=\pm(e,f)+ (g,h)+(n_1,n_2)$, then $-2=\pm e+g+n_1$. So, Lemma \ref{5} gives that $R_{\alpha_1}$ is strongly 2-nil-clean -- a contradiction. If $(-2,2)=(e,f)\pm (g,h)+(n_1,n_2)$, then via a similar argument $R_{\alpha_1}$ is strongly 2-nil-clean -- again a contradiction. Finally, if $(-2,2)=-(e,f)- (g,h)+(n_1,n_2)$, then $2=-f-h+n_2$ and so $-2=f+h-n_2$. Applying Lemma \ref{5}, it must be that $R_{\alpha_2}$ is strongly 2-nil-clean, which is an absurd. Thus, $(-2,2)$ can not express as weakly strongly 2-nil-clean decomposition. So, either $\alpha_1$ or $\alpha_2 $ has to be strongly 2-nil-clean, as stated.

Conversely, assume that $R_{{\alpha}_1}$ is a weakly strongly 2-nil-clean ring that is {\it not} strongly 2-nil-clean and $R_{\alpha}$ is strongly 2-nil-clean for all other $\alpha\neq \alpha_1$. We can now write $\prod_{\alpha\neq \alpha_1} R_{\alpha}\simeq R_1\times R_2$, where $3\in nil(R_1)$ and $2\in nil(R_2)$. So, we arrive at the isomorphism $\prod_{\alpha} R_{\alpha}\cong R_{\alpha_1}\times R_1\times R_2$, where, as proven above in Lemmas~\ref{3} and \ref{2}, $R_{\alpha_1}$ is weakly strongly 2-nil-clean that is {\it not} strongly 2-nil-clean, $R_1$ is a strongly 2-nil-clean ring and $R_2$ is a strongly nil-clean ring. Now, let $(r,r_1,r_2)\in\prod R_{\alpha}$. Then, there are $e,f\in Id(R)$ and $n\in nil(R)$ all commuting such that $r=\pm e\pm f+n$. We, hereafter, differ three cases as follows:

\medskip

\noindent{\bf Case (1):} If $r=e+f+n$, then Lemma \ref{2nil} enables us that $r_1=e_1+f_1+n_1$ for some $e_1,f_1\in Id(R_1)$ and $n_1\in nil(R_1)$ which all commute, and that there exists $e_2\in Id(R_2)$ and $n_2\in nil(R_2)$ which commute such that $r_2=e_2+n_2$. It follows now that $$(r,r_1,r_2)=(e,e_1,e_2)+(f,f_1,0)+(n,n_1,n_2),$$ and, surely, $(e,e_1,e_2),(f,f_1,0)\in Id(\prod_{\alpha} R_{\alpha})$ and $(n,n_1,n_2)\in nil(\prod_{\alpha} R_{\alpha})$ that commute with each other.

\medskip

\noindent{\bf Case (2):} If $r=e-f+n$, then Lemma \ref{2nil} tells us that $r_1=e_1-f_1+n_1$ for some $e_1,f_1\in Id(R_1)$ and $n_1\in nil(R_1)$ which all commute, and that there exists $e_2\in Id(R_2)$ and $n_2\in nil(R_2)$ which commute such that $r_2=e_2+n_2$. It follows now that $$(r,r_1,r_2)=(e,e_1,e_2)-(f,f_1,0)+(n,n_1,n_2),$$ and, evidently, $(e,e_1,e_2),(f,f_1,0)\in Id(\prod_{\alpha} R_{\alpha})$ and $(n,n_1,n_2)\in nil(\prod_{\alpha} R_{\alpha})$ that commute with each other.

\medskip

\noindent{\bf Case (3):} If $r=-e-f+n$, then Lemma \ref{2nil} teaches us that $r_1=-e_1-f_1+n_1$ for some $e_1,f_1\in Id(R_1)$ and $n_1\in nil(R_1)$ which all commute, and that there exists $e_2\in Id(R_2)$ and $n_2\in nil(R_2)$ which commute such that $r_2=e_2+n_2$, whence $r_2=-e_2+n_2+2e_2$ and $n_2+2e_2\in nil(R_2)$. It follows now that $$(r,r_1,r_2)=-(e,e_1,e_2)-(f,f_1,0)+(n,n_1,n_2+2e_2),$$ and, easily, $(e,e_1,e_2),(f,f_1,0)\in Id(\prod_{\alpha} R_{\alpha})$ and $(n,n_1,n_2+2e_2)\in nil(\prod_{\alpha} R_{\alpha})$ that commute with each other, thus finishing the entire proof.
\end{proof}

\end{document}
